\documentclass[11pt]{article}
\usepackage[utf8]{inputenc}
\usepackage{amsmath,amssymb}
\usepackage[margin=1in]{geometry}
\usepackage{hyperref}
\emergencystretch=3em

\title{The boundary tail of the standard integral is exponentially small}
\author{Claude, with Daniel Murfet}
\date{2026-09-06}

\begin{document}
\maketitle
\sloppy

\begin{abstract}
Third \S4.2 foothold: the one-dimensional population version of the remainder estimate
\texttt{eq:Deltaupper}. The paper's standard integral lives on the chart $[0,b]$; the
fluctuation-function identity of unit~22 lives on $(0,\infty)$. We bound the difference ---
the tail over $(b,\infty)$ --- by $e^{\beta a^2/2}e^{-(\beta/4)nb^{2k}}$ times an explicit Gamma
constant, for $n\ge1$, and conclude that the chart integral equals
$(2k)^{-1}n^{-(h+1)/(2k)}S_{(h+1)/(2k)}(a)$ up to $O(e^{-\varepsilon n})$ with
$\varepsilon=(\beta/4)b^{2k}$. Zero \texttt{sorry}/\texttt{axiom}.
\end{abstract}

\section{Setting}
In the proof of \texttt{thm:TaylorTree} the state-density integral is cut at $T=nb^{2|k|}$ and the
remainder $\Delta$ is bounded by $e^{-\beta T/4}$ times a convergent multi-sum. Stripped of the tree
combinatorics, the analytic content is: on $u\ge b$ the exponent $-\beta nu^{2k}+\beta\sqrt n u^k a$
is at most $-\tfrac\beta4 nb^{2k}-\tfrac\beta4u^{2k}+\tfrac\beta2a^2$, so the tail is dominated by a
fixed Gamma integral times $e^{-(\beta/4)nb^{2k}}$.

\section{The things formalised}
{\small
\begin{verbatim}
theorem standardIntegrand_le ... :
    Real.exp (-β*n*u^(2*k) + β*Real.sqrt n*u^k*a)
      ≤ Real.exp (β*a^2/2) * Real.exp (-(β/2*n) * u^(2*k))
theorem standardIntegrand_integrableOn ... :
    IntegrableOn (fun u => u^h * exp(...)) (Ioi 0)
theorem standardIntegral1D_tail_le (β n a b : ℝ) (h k : ℕ) (hβ : 0 < β) (hn : 1 ≤ n)
    (hb : 0 < b) (hk : 0 < k) :
    (∫ u in Ioi b, u^h * exp(...))
      ≤ exp(β*a^2/2) * exp(-(β/4)*n*b^(2*k))
        * ((β/4)^(-(((h:ℝ)+1)/(2*k))) * (1/(2*(k:ℝ))) * Gamma(((h:ℝ)+1)/(2*k)))
theorem standardIntegral1D_bounded_approx ... :
    |(∫ u in Ioc 0 b, u^h * exp(...))
        - (1/(2*(k:ℝ))) * n^(-(((h:ℝ)+1)/(2*k))) * fluctuation β (((h:ℝ)+1)/(2*k)) a|
      ≤ (the same bound)
\end{verbatim}
}

\section{Proof strategy}
\texttt{standardIntegrand\_le} is the Gaussian domination of unit~23 evaluated at $t=nu^{2k}$, using
$\sqrt{nu^{2k}}=\sqrt n\,u^k$. Integrability on $(0,\infty)$ follows by \texttt{Integrable.mono'}
against $e^{\beta a^2/2}u^h e^{-(\beta/2)nu^{2k}}$, which is
\texttt{integrableOn\_rpow\_mul\_exp\_neg\_mul\_rpow} after \texttt{rpow\_natCast}. For the tail, on
$u>b$ we have $u^{2k}\ge b^{2k}$ and $n\ge1$, so
$-\tfrac\beta2nu^{2k}\le-\tfrac\beta4nb^{2k}-\tfrac\beta4u^{2k}$ (\texttt{nlinarith} with the two
product hints $\beta n(u^{2k}-b^{2k})\ge0$ and $\beta(n-1)u^{2k}\ge0$); the tail integral is then
bounded by \texttt{setIntegral\_mono\_on} against $C\cdot u^he^{-(\beta/4)u^{2k}}$, enlarged from
$(b,\infty)$ to $(0,\infty)$ by \texttt{setIntegral\_mono\_set}, and evaluated by
\texttt{integral\_rpow\_mul\_exp\_neg\_mul\_rpow}. The chart approximation splits $(0,\infty)$ as
$(0,b]\cup(b,\infty)$ (\texttt{Ioc\_union\_Ioi\_eq\_Ioi}, \texttt{setIntegral\_union}) and invokes
unit~22 for the full-line value.

\section{Roadblocks and resolutions}
\begin{itemize}
\item \texttt{Integrable.const\_mul} returns a bare \texttt{Integrable} (an \texttt{And}), so
dot-notation \texttt{.mono\_set} fails with ``\texttt{And.mono\_set}''; bind it first with the
\texttt{IntegrableOn} type ascribed.
\item \texttt{pow\_le\_pow\_left} is now \texttt{pow\_le\_pow\_left₀} on this pin.
\item The a.e.\ nonnegativity hypothesis of \texttt{setIntegral\_mono\_set} is stated with
\texttt{≤ᵐ[μ.restrict t]}; unfold \texttt{Filter.EventuallyLE} then \texttt{ae\_restrict\_iff'} to
reduce it to a pointwise statement on the set, and use \texttt{HasSubset.Subset.eventuallyLE} for the
set inclusion.
\item The Gamma-lemma output carries $(-(h+1))/(2k)$; \texttt{push\_cast} then \texttt{neg\_div}
aligns it with the stated $-((h+1)/(2k))$.
\end{itemize}

\section{What was Mathlib, what was new}
Mathlib: the Gamma-integral pair, \texttt{Integrable.mono'}, \texttt{setIntegral\_mono\_on},
\texttt{setIntegral\_mono\_set}, \texttt{setIntegral\_union}, \texttt{Ioc\_union\_Ioi\_eq\_Ioi}.
New: integrability of the standard integrand, the explicit exponential tail bound with its Gamma
constant, and the chart-vs-half-line approximation --- the 1D skeleton of $\Delta=O(e^{-\varepsilon n})$.

\section{Lessons}
A domination lemma factored out once (unit~23) pays twice: it gives integrability and the tail bound
with the same majorant. State tail bounds with the explicit constant rather than an
$O(\cdot)$ --- the constant is one Gamma evaluation away and makes the result quotable.

\section{Follow-ups}
The 1D population Taylor-tree pipeline now has all its analytic steps (kernel identity, insertions,
tail). What remains multivariate: the tree-coefficient series interchange and the state-density
Mellin machinery.
\end{document}
